\section{Objetivos e hipótesis}

\subsection{Objetivo general}
Evaluar el desempeño de configuraciones de una red neuronal convolucional ligera
para el reconocimiento óptico de Braille Grado 1 en dispositivos móviles, con el
fin de identificar la que ofrece el mejor compromiso entre precisión (incluida la
robustez), eficiencia y oportunidad de la retroalimentación docente.

\subsection{Objetivos específicos}
\begin{enumerate}
  \item Comparar la precisión y la robustez de las configuraciones evaluadas
  mediante mAP@0.5, F1 macro, Kappa de Cohen y AUC-ROC macro.
  \item Determinar el efecto de la cuantización INT8 sobre la latencia, la
  memoria y la precisión del modelo.
  \item Evaluar el efecto del sistema en el tiempo de revisión por página de los
  docentes.
\end{enumerate}

\subsection{Hipótesis}
La hipótesis general anticipa que YOLOv11 cuantizada a INT8 ofrece el mejor
compromiso entre precisión y eficiencia, y reduce el tiempo de revisión docente.

\begin{itemize}
  \item \textbf{H1.} YOLOv11 supera a YOLOv8n en mAP@0.5 y F1 macro (Wilcoxon;
  $p < .05$) y mantiene la ventaja en condiciones adversas.
  \item \textbf{H2.} INT8 reduce la latencia y la memoria con una degradación de
  mAP@0.5 no superior al umbral que se fije en la corrida piloto.
  \item \textbf{H3.} El sistema reduce el tiempo de revisión por página respecto
  del procedimiento habitual (Wilcoxon; $p < .05$).
\end{itemize}
